\unnumchapter{Выводы}

По каждой из поставленных задач получены следующие результаты.

\begin{enumerate}
  \item Проведён обзор средств управления федерацией, наблюдаемости,
        аварийного восстановления и контроля политик. Установлено, что
        существующие федеративные планировщики выбирают кластер по
        статическим атрибутам, а ни одно из рассмотренных решений не
        объединяет политики, текущие показатели качества и состояние
        кластеров в одно решение о размещении.
  \item Развёрнут экспериментальный стенд из управляющего кластера и двух
        кластеров федерации с тестовым сервисом, шлюзом и генератором
        нагрузки. Стенд и каждый эксперимент воспроизводятся из
        репозитория одной командой.
  \item Построен слой наблюдаемости федерации: метрики кластеров
        агрегируются через Thanos, определены локальные и глобальные
        показатели уровня обслуживания, реализована сквозная трассировка с
        разложением задержки по звеньям. Показано, что доступность
        федерации необходимо измерять в точке входа трафика: сумма
        серверных показателей кластеров не видит отказа кластера.
  \item Разработан каталог кластеров, хранящий атрибуты кластеров, их
        текущее состояние, журнал решений о размещении со всеми
        кандидатами, журнал нарушений политик и инцидентов. Метрики
        каталога совпадают с исходными рядами в 98--99\% снимков.
  \item Реализованы двухуровневые политики размещения. В режиме контроля
        не размещён ни один из 15 нарушителей, все попытки зафиксированы в
        журнале в течение не более 10~с.
  \item Разработан планировщик уровня федерации, оценивающий кластеры по
        текущим доле ошибок, задержке и загрузке, и наивный планировщик,
        выбирающий кластеры по очереди среди тех же допустимых кандидатов.
  \item Разработан контроллер аварийного восстановления на основе конечного
        автомата с асимметричными переходами. Время обнаружения отказа
        составило 10.6--12.8~с, время восстановления~--- 40.7--42.9~с при
        целевых значениях 60~с и 5~мин.
  \item Проведены эксперименты по трём сценариям отказов. Контроллер
        восстановления снижает число потерянных запросов при отказе
        кластера на 87.1\%. Разрабатываемый планировщик снижает долю ошибок
        при деградации сети на 45.6\% относительно наивного, превосходя его
        в каждой паре прогонов.
  \item Проведён анализ результатов. Преимущество планировщика сохраняется
        в 64 из 65 комбинаций весов и исчезает, когда оценка не учитывает
        задержку. Обнаружено, что проверки работоспособности с таймаутом,
        не учитывающим задержку сети, превращают сетевую деградацию в отказ
        сервиса.
\end{enumerate}
